\documentclass[10pt,a4paper]{article}
\usepackage[a4paper,top=18mm,bottom=19mm,left=19mm,right=19mm,headheight=13.6pt]{geometry}
\usepackage{fontspec}
\IfFontExistsTF{TeX Gyre Termes}
  {\setmainfont{TeX Gyre Termes}}{\setmainfont{Liberation Serif}}
\IfFontExistsTF{TeX Gyre Heros}
  {\setsansfont{TeX Gyre Heros}}{\setsansfont{Liberation Sans}}
\usepackage{amsmath,amssymb}
\usepackage{graphicx}
\usepackage{xcolor}
\usepackage{hyperref}
\usepackage{bookmark}
\usepackage{enumitem}
\usepackage{titlesec}
\usepackage{fancyhdr}
\usepackage{microtype}
\usepackage{float}
\definecolor{AIDRBlue}{HTML}{18364A}
\definecolor{AIDRLink}{HTML}{245F7A}
\hypersetup{colorlinks=true,linkcolor=AIDRLink,urlcolor=AIDRLink,citecolor=AIDRLink,
pdfauthor={AIDRBench authors},
pdftitle={Reliable demand-response commitments from AI data centres under service constraints}}
\urlstyle{same}
\setlength{\parindent}{0pt}
\setlength{\parskip}{5.5pt plus 1pt minus 0.5pt}
\setlength{\emergencystretch}{2em}
\widowpenalty=10000\clubpenalty=10000\displaywidowpenalty=10000
\setlist[enumerate]{leftmargin=*,itemsep=4pt,topsep=4pt}
\setlist[itemize]{leftmargin=*,itemsep=2pt,topsep=3pt}
\titleformat{\section}{\Large\sffamily\bfseries\color{AIDRBlue}}{}{0pt}{}
\titleformat{\subsection}{\large\sffamily\bfseries\color{AIDRBlue}}{}{0pt}{}
\titlespacing*{\section}{0pt}{17pt}{6pt}
\titlespacing*{\subsection}{0pt}{12pt}{4pt}
\pagestyle{fancy}
\fancyhf{}
\fancyhead[L]{\small\sffamily AI data centres and demand response}
\fancyhead[R]{}
\fancyfoot[C]{\small\thepage}
\makeatletter
\setlength{\@fptop}{0pt}
\setlength{\@fpsep}{12pt}
\setlength{\@fpbot}{0pt plus 1fil}
\makeatother
\renewcommand{\topfraction}{0.95}
\renewcommand{\bottomfraction}{0.95}
\renewcommand{\textfraction}{0.05}
\renewcommand{\floatpagefraction}{0.80}
\title{\sffamily\bfseries Reliable demand-response commitments from AI data centres under service constraints}
\author{\sffamily [AUTHOR NAMES AND AFFILIATIONS]}
\date{}
\begin{document}
\maketitle
\thispagestyle{fancy}




\emph{Correspondence: [CORRESPONDING AUTHOR EMAIL]}


\section*{Abstract}

The expansion of artificial intelligence intensifies a central challenge of the energy transition: meeting rising electricity demand with reliable, low-carbon supply. Data-centre demand response could ease these pressures, but planners need reliable commitments that preserve computing service through repeated calls. Here we combine production records, power measurements and scheduling simulations in a model of 576 graphics processing units. At 10\% deferrable work, retaining task sizes and durations instead of counting jobs reduces the independently confirmed eight-hour offer from 2.95 to 1.47 kW within the tested eight-week distribution. Checks of available load and task deadlines distinguish infeasible requests from failures that improved control could address. For separately qualified reference products, four- and eight-hour participation thresholds are US\$183 and US\$330 per committed kilowatt-year under illustrative waiting costs and zero fixed fees. These results connect capacity qualification to service selection, helping operators avoid infeasible commitments and compare compensation for reliable grid support.


\section*{Introduction}

The global energy transition requires electricity systems to cut emissions while meeting rising demand from electrification and artificial intelligence (AI).\textsuperscript{1} Data centres intensify this challenge where computing demand is concentrated and grows faster than power infrastructure can be built.\textsuperscript{1,2} In these regions, limited electricity supply can constrain AI expansion, while grid planners must keep power reliable and affordable.\textsuperscript{2,3} Shifting some computation to other times offers a way to make better use of existing infrastructure.\textsuperscript{3} To support grid planning, data centres must repeatedly reduce power when requested, complete delayed tasks on time and keep costs acceptable. The reductions they can reliably promise therefore help determine their contribution to the energy transition.


A field demonstration on 256 graphics processing units (GPUs) reduced power by 25\% for three hours while meeting the tested service requirements.\textsuperscript{3} Previous studies have shifted tasks between times or locations and controlled server power to reduce peaks and support demand response.\textsuperscript{4–8} Related work schedules computing with cleaner electricity,\textsuperscript{9–12} examines task pausing and GPU power limits,\textsuperscript{13,14} and considers advance notice and grid connection.\textsuperscript{15–18} A recent preprint relates reliable power reductions to the participating workload, response duration, reliability and the need to complete postponed work.\textsuperscript{19} These advances raise a practical question: can simplifying task sizes and timing lead an operator to promise power reductions it cannot deliver?


Some computing must answer users immediately; some model training and predictions prepared for later use can wait within agreed deadlines. Production records distinguish these task types and their priorities,\textsuperscript{20} but operators must decide which work may actually be delayed. At a fixed amount of deferrable work, hourly demand still depends on how many GPUs each task needs and how long it runs. Delaying tasks during a response also leaves work that must be completed before its deadline.\textsuperscript{6,13,21,22} A centre can therefore have enough work over a week yet too little reducible load during a call or too little time afterwards.


Whether a centre can deliver a response and whether participation pays are separate questions. On a shared local electricity network, computing demand interacts with solar generation and battery storage.\textsuperscript{23–27} Equipment that allows more solar generation to connect may not increase the power reduction available during grid calls. Delaying work also creates waiting costs, while missed deadlines and programme access can add expenses. We therefore first establish which commitments can be delivered and then compare the costs of the services that pass.


Here we identify how omitting task-size information can cause an operator to overstate deliverable response. AIDRBench combines computing tasks modelled from production records, power measurements from four GPUs and a simulated local electricity network. We compare two descriptions of the same weekly work: job counts alone, or counts adjusted for requested GPUs and task duration.\textsuperscript{28} We select capacities using one set of simulations and test them on a separate set. We then diagnose failures and compare participation costs for separately tested reference services. Building on previous capacity assessments,\textsuperscript{19} we show when operators should reduce a request, reconsider scheduling or compare compensation among deliverable services.


\section*{Results}

\subsection*{Workload composition limits which computing can support a commitment}

We first examined what fraction of computing could plausibly enter the deferrable pool. Across 40.52 million execution spans in the released Alibaba summary,\textsuperscript{20} online inference accounted for 54.5\% of requested GPU-hours, offline inference for 21.5\% and training for 19.4\% (Fig. 1b). Low-priority training and offline inference together accounted for 22.0\%. These descriptive resource-time shares identify a candidate pool from which an operator might authorise a smaller amount of work to wait.


The primary scenarios allowed 5\%, 10\% or 20\% of total offered work to wait, preserving the observed class mix and varying permission within low-priority batch work. The 40\% comparison required almost all training and offline inference to opt in. The 60\% case additionally reassigned 19.1 percentage points of total work from online to offline inference (Supplementary Table 2). Online jobs remained ineligible; these larger fractions therefore describe additional business assumptions.


Each configuration retained 576 GPUs and 374.4 GPU-h of offered work per hour. Assigning approximately the same fraction of GPUs as eligible work kept mean utilisation near 65\% in both pools. The primary comparison therefore changes permission together with supporting allocation; controls at 10\% eligibility isolate allocation (Fig. 5). Four-GPU board-power calibration constrained the batch-work conversion, whereas rigid workloads used an engineering power proxy (Supplementary Fig. 2; Supplementary Tables 1–3).


\begin{figure}[htbp]
\centering
\begin{minipage}{\textwidth}
\setlength{\parskip}{0pt}
\pdfbookmark[1]{Figure 1 | AI workload composition, deferred scheduling and assessment of demand-response commitments}{main-figure-1}
{\fontsize{10}{12}\selectfont\sffamily\bfseries\color{AIDRBlue} Figure 1 | AI workload composition, deferred scheduling and assessment of demand-response commitments\par}\vspace{4pt}
\centering
\makebox[\textwidth][c]{\includegraphics[width=180mm]{\detokenize{figures/AIDRBench_Figure_1.pdf}}}
\end{minipage}
\end{figure}

\subsection*{Smaller eligible pools support smaller independently tested offers}

Development-selected single-event offers of 2.95, 5.90 and 11.80 kW qualified on 300 independent scenarios at 5\%, 10\% and 20\% eligibility, respectively, corresponding to 1.6–5.9\% of operating peak demand (Fig. 2a; Supplementary Table 4). The 40\% and 60\% comparisons qualified at 22.19 and 34.00 kW under their additional assumptions. The proportional low-share response follows the shared job-template scaling and allocation policy. Relaxed perfect-information (PI) tolerance statistics provide a planning comparison but do not guarantee individual work-group feasibility (Methods).


The same selected power qualified at four and eight hours on the tested grid, with success decreasing from 295/300 to 292/300 in every configuration; one-sided 95\% probability lower bounds were 0.9662 and 0.9533. Zero-, two- and six-hour notice gave identical paired success flags (Fig. 2c). This result belongs to the specified eager baseline, release constraints and controller; the baseline-dependent supply bound explains why notice cannot repair hours with insufficient reducible load.


\begin{figure}[htbp]
\centering
\begin{minipage}{\textwidth}
\setlength{\parskip}{0pt}
\pdfbookmark[1]{Figure 2 | Single-event response capacity and reliability across workload eligibility, event duration and advance notice}{main-figure-2}
{\fontsize{10}{12}\selectfont\sffamily\bfseries\color{AIDRBlue} Figure 2 | Single-event response capacity and reliability across workload eligibility, event duration and advance notice\par}\vspace{4pt}
\centering
\makebox[\textwidth][c]{\includegraphics[width=180mm]{\detokenize{figures/AIDRBench_Figure_2.pdf}}}
\end{minipage}
\end{figure}

\subsection*{Workload timing determines whether a repeated commitment transfers}

At 10\% work eligibility, we tested four calls using a controller that enforces every overlapping recovery obligation. In the reference configuration—65\% offered utilisation, reference deadlines and a 10\% GPU pool—the four- and eight-hour offers were 5.60 and 4.42 kW, both with 16-h start spacing. When successful series required zero missed work, independent confirmation gave 297/300 and 296/300 successes, with one-sided 95\% probability lower bounds of 0.9752 and 0.9706 (Supplementary Fig. 9a). Both the 1\% and zero-miss standards selected these capacities on the local grid (Supplementary Table 18). Their agreement provides no basis for applying a uniform capacity discount solely because the success criterion becomes stricter.


Even half the reference single-event offer could fail after a change in workload representation. At 10\% eligibility, a 2.95-kW request—50\% of the 5.898243-kW single-event offer—passed all ten simulations in each of eight observed weeks under submission-count weighting. Under task-resource-time weighting, only two weeks passed all ten; two failed all ten and four had mixed outcomes (Fig. 3d,e). The descriptive total fell from 80/80 to 31/80. Weekly class totals, synthetic service rules, community profiles and event clocks were paired, and every no-response baseline met zero-miss service. This 2.95-kW stress request concerns the 10\% configuration; the numerically equal single-event offer at 5\% eligibility is a separate product.


Permuting whole hourly blocks separated workload weighting from the effect of hourly order and alignment. Count-weighted success remained 80/80, whereas resource-time success rose to 55/80 at the 1\% standard and 54/80 at zero misses (Fig. 3d). Hourly rearrangement improved this comparison but did not restore count-weighted performance. It changes alignment with both calls and community demand, so the improvement cannot be assigned to autocorrelation alone. At the larger 4.42-kW request, all 80 chronological resource-time realisations failed; full cross-scores are retained in Supplementary Table 20.


Reducing the request restored qualification under the resource-time description. We selected offers on 100 new simulations and froze them before testing 300 further simulations, each independently drawing one of the eight fixed weekly profiles with equal probability. The count-weighted description selected 2.9491 kW and the resource-time description 1.4746 kW; each achieved 300/300 successes under both service standards, with a one-sided 95\% lower bound of 0.9911 (Fig. 3a; Supplementary Fig. 9; Supplementary Table 22). At the unchanged 2.9491-kW comparator, resource-time success was 120/300, with all 180 failures encountering supply shortage. Thus the simplified description supported a tested offer twice as large. This ratio compares finite-grid offers conditional on the eight-week empirical distribution; it does not estimate a universal or continuous optimal-capacity ratio.


\begin{figure}[htbp]
\centering
\begin{minipage}{\textwidth}
\setlength{\parskip}{0pt}
\pdfbookmark[1]{Figure 3 | Repeated-response offers and delivery outcomes under submission-count weighting, task-resource-time weighting and hourly reordering}{main-figure-3}
{\fontsize{10}{12}\selectfont\sffamily\bfseries\color{AIDRBlue} Figure 3 | Repeated-response offers and delivery outcomes under submission-count weighting, task-resource-time weighting and hourly reordering\par}\vspace{4pt}
\centering
\makebox[\textwidth][c]{\includegraphics[width=180mm]{\detokenize{figures/AIDRBench_Figure_3.pdf}}}
\end{minipage}
\end{figure}

\subsection*{Supply and feasibility checks identify what a failed request requires}

The matched baseline makes a simple necessary supply condition available. In each event hour, baseline power above community and idle-plus-rigid demand must reach 95\% of the request. A deficit makes that hour undeliverable even with all flexible execution stopped. This condition accounts for all 49 chronological resource-time failures at 2.95 kW and 25 failures after permutation (Fig. 4a,b). Full queue replay establishes that supply, rather than an additional recovery or scheduling violation, accounts for every failure in the original-order comparison. The screen is not sufficient in general: one of the 55 permuted cases without shortage still failed the zero-miss endpoint. Strict-window and causal tests address these remaining constraints.


Strict task-window feasibility then distinguished failures beyond that supply condition. At a fixed 5.90-kW request, a reference scenario that passed the supply screen failed under causal control but admitted a full-information schedule meeting all electrical obligations with zero missed work (Fig. 3c; Supplementary Table 19). Halving its deadline slack made the request infeasible under the same delivery and recovery criteria; doubling GPU allocation did not restore feasibility. Each variant used its own matched baseline. The feasible reference case leaves room for controller improvement, whereas proven infeasibility calls for changing the request or operating terms. Neither finding follows from the causal failure alone. The broader structural confirmation separates joint success from overlapping supply and deadline failures (Supplementary Fig. 6).


These checks define different decisions (Fig. 1c): revise requests that exceed supply or strict-window feasibility, investigate control where a feasible PI schedule exists, and qualify the causal policy independently before comparing prices. Qualification uses the complete success and failure sample. In the matched same-clock control, removing preceding calls left final-call electrical outcomes unchanged (Supplementary Table 15).


At the qualified reference offers, the full response–recovery trajectory illustrates why an event cannot be assessed in isolation: deferred work accumulates during the calls and is processed afterwards (Fig. 4c,d). This seed-selected example complements, rather than replaces, the independent qualification counts.


\begin{figure}[htbp]
\centering
\begin{minipage}{\textwidth}
\setlength{\parskip}{0pt}
\pdfbookmark[1]{Figure 4 | Available-load shortages during repeated calls and power and backlog trajectories through recovery}{main-figure-4}
{\fontsize{10}{12}\selectfont\sffamily\bfseries\color{AIDRBlue} Figure 4 | Available-load shortages during repeated calls and power and backlog trajectories through recovery\par}\vspace{4pt}
\centering
\makebox[\textwidth][c]{\includegraphics[width=180mm]{\detokenize{figures/AIDRBench_Figure_4.pdf}}}
\end{minipage}
\end{figure}

\subsection*{Allocation changes system value without necessarily increasing the offer}

At fixed 10\% work eligibility, increasing flexible GPU allocation from 10\% to 20\% or 30\% left the tested 5.90-kW single-event offer at 295/300 four-hour and 292/300 eight-hour successes. The corresponding relaxed-PI tolerance statistics were also unchanged (Fig. 5a; Supplementary Fig. 4c). Reducing the rigid-class active-power proxy from 300 to 225 or 150 W per GPU preserved these baseline-relative responses but changed operating peak demand from 193.44 to 173.53 or 153.63 kW (Fig. 5b). Thus, eligible arrivals constrained this response, while rigid power remained relevant to facility demand and proportional cost accounting (Fig. 5c,d).


The same allocation could matter more for integrating photovoltaic (PV) generation, assessed with and without a battery energy storage system (BESS). At 10\% eligibility, raising flexible GPU allocation from 10\% to 20\% increased the no-BESS PV hosting gain from 5.40 to 24.99 kW (Fig. 5e; Supplementary Fig. 4d). Strict repeat solves retained these gains. Utilisation of a fixed 500-kW PV system improved by only 0.0057 percentage points without BESS (95\% paired interval, 0.0019–0.0106). This small LP result was unchanged by tighter feasibility tolerances, whereas the corresponding BESS gain was numerically zero (Supplementary Fig. 10c,d; Supplementary Table 24). Hosting and utilisation therefore measure different benefits; absolute hosting and energy outcomes appear in Supplementary Fig. 7.


\begin{figure}[htbp]
\centering
\begin{minipage}{\textwidth}
\setlength{\parskip}{0pt}
\pdfbookmark[1]{Figure 5 | Response planning, participation costs and solar hosting across GPU allocations and rigid-load power assumptions}{main-figure-5}
{\fontsize{10}{12}\selectfont\sffamily\bfseries\color{AIDRBlue} Figure 5 | Response planning, participation costs and solar hosting across GPU allocations and rigid-load power assumptions\par}\vspace{4pt}
\centering
\makebox[\textwidth][c]{\includegraphics[width=180mm]{\detokenize{figures/AIDRBench_Figure_5.pdf}}}
\end{minipage}
\end{figure}

\subsection*{Access fees and operating costs affect different participation decisions}

Once a duration product qualified, its operating exposure and access charges could be valued separately. In the 10\% eligibility single-event case, the assumed US\$25,000 annual site fee accounted for about 95\% of the US\$863.26 per offered kW-year threshold (Supplementary Table 17). To distinguish fee allocation from operating consequences, we compared complete response-and-recovery ledgers for the qualified repeated products. All amounts use proportional 1-MW accounting and retain failed confirmation trajectories.


For the four- and eight-hour products qualified with a zero-miss success criterion, mean additional waiting exposure per four-call series was 16,491 and 24,509 GPU-h×h, respectively (Fig. 6a; Supplementary Table 21). At the illustrative waiting price of US\$0.005/GPU-h/h and twelve independent series per year, the no-fixed-fee thresholds were US\$183.30 and 330.11 per offered accounting kW-year (Fig. 6b,c). Their capacity difference and waiting exposure explain this conditional comparison; the valuations have not been measured as an operator’s willingness to accept.


At equal illustrative capacity prices of US\$250 per kW-year and zero site fee, the four-hour product’s fifth-percentile annual net value was US\$1,932, compared with −US\$1,832 for eight hours. Conditional on participation, equal value requires an eight-hour price of 1.267 times the four-hour price plus US\$97.93 per kW-year; the full boundary also requires non-negative net value (Fig. 6d). The coefficient is the ratio of the selected capacities. When waiting has no incremental price, the equal-price ranking reverses below approximately US\$25.57 per kW-year (Supplementary Fig. 5). A shared fee can change whether either option beats non-participation, but a common fee cancels between participating products.


The price comparison applies to the qualified reference-distribution products. The lower resource-time offer requires its own operating ledger before those prices can be used for that workload. Waiting and missed-work valuations are varied explicitly, with single-event hardware and monetary sensitivity in Supplementary Fig. 8. Contract penalties, application checkpoint costs and hardware ageing remain unpriced, and annual accounting combines independent series.


\begin{figure}[htbp]
\centering
\begin{minipage}{\textwidth}
\setlength{\parskip}{0pt}
\pdfbookmark[1]{Figure 6 | Additional waiting, participation costs and service choice for qualified four- and eight-hour repeated responses}{main-figure-6}
{\fontsize{10}{12}\selectfont\sffamily\bfseries\color{AIDRBlue} Figure 6 | Additional waiting, participation costs and service choice for qualified four- and eight-hour repeated responses\par}\vspace{4pt}
\centering
\makebox[\textwidth][c]{\includegraphics[width=180mm]{\detokenize{figures/AIDRBench_Figure_6.pdf}}}
\end{minipage}
\end{figure}

\section*{Discussion}

The main finding is a quantitative commitment error: removing task resource requirements from an otherwise matched workload description supported a tested offer twice as large. Both offers subsequently passed independent confirmation, so the resource-time result specifies a usable lower commitment within the model. Weekly totals alone concealed the low-supply hours that invalidated the larger request. Paired permutation improved resource-time delivery without restoring count-weighted performance, separating workload weights from their hourly arrangement and alignment. These results extend availability-based contract descriptions\textsuperscript{19} by testing the request through explicit task queues and repeated recovery, then separating physical impossibility from a controller's failure to realise a feasible schedule.


The diagnosis also determines what can usefully change. An instantaneous deficit rules out successful delivery in the affected hour. If that condition passes, strict task-window infeasibility and a feasible PI schedule with causal failure require different responses: revise operating terms in the first case, investigate control and information limits in the second. These scenario-level findings must then be assessed through the declared probabilistic qualification rule. More GPUs, a stricter service label or a low estimated price cannot substitute for that sequence of evidence.


These findings depend on what a contract counts as delivery. Our counterfactual baseline immediately processes available work and retains the same rigid demand and powered GPU idle floor as the response trajectory. It therefore fixes each hour's maximum baseline-relative reduction. Notice cannot lift this ceiling, although it may improve service or recovery scheduling when supply is sufficient. Historical customer-baseline load or a fixed import ceiling defines a different obligation and can change the value of pre-positioning work. PJM's capacity framework distinguishes guaranteed load drop from a firm service level.\textsuperscript{29} Our fixed-duration repeated reduction is closest in intent to a guaranteed-reduction commitment; its matched counterfactual baseline and recovery rules are explicit modelling choices, rather than an implementation of PJM settlement or baseline estimation.


Economic assessment follows qualification and exposes quantities before assigning prices. Waiting and missed-work exposure remain inspectable independently of their valuations. The preferred duration can change when waiting is free, whereas a common access fee affects entry rather than the ranking of participating products. Comparing annual net values with non-participation avoids confusing a low unit-cost threshold with the most valuable contract. The capacity ratio and stated valuations determine the price boundary, whose use remains conditional on the workload for which delivery was qualified.


For planners and operators, the useful output is a conditional commitment with a declared workload, baseline, service standard, capacity and call arrangement. Its delivery evidence and operating ledger must be re-evaluated when the timing structure changes. This is consistent with capacity assessment for other resources whose decisions are coupled over time.\textsuperscript{30,31} A larger nominal flexible share, more recovery hardware or cheaper access cannot substitute for that assessment.


The study represents a 576-GPU installation with 58 GPUs assigned to eligible work in the primary case. Its 193.44-kW operating peak includes 61.47 kW of node overhead and GPU idle demand, limiting the removable share before task timing is considered. This denominator and the 10\% work permission distinguish the scenario from demonstrations with different participating workloads and hardware.\textsuperscript{3} The trace supplies allocated task resource-time, while deadlines, permissions, divisible execution and recovery behaviour remain model assumptions. Hardware evidence calibrates board power rather than application pause–restart behaviour. Ignoring checkpoint overhead and indivisibility enlarges the scheduling possibilities at a fixed physical baseline; a causal controller's observed failure rate still cannot be treated as a lower bound on real-system failure. The eight-week empirical distribution, pointwise qualification and finite grids delimit the capacity comparison, and independent annual series delimit economic extrapolation.


\section*{Methods}

\subsection*{Study design and capacity definition}

The primary comparison followed five declared workload-eligibility scenarios through perfect-information (PI) planning, independent causal-controller testing, repeated programmes, PV integration and participation costs. Eligibility was the share of offered GPU-hours permitted to wait. Capacity was evaluated against the same delivery and service criteria within each declared call programme.


Let Z describe an episode's initial queue, community load, job arrivals and deadlines, hardware and event time, drawn from distribution \(\mathcal D\). For fixed policy \(\pi\), reduction R, duration H and notice N, \(I_\pi(R,H,N;Z)=1\) if all delivery and service criteria are met. Firm capacity was defined as


\begin{equation*}
F_q^{\pi}(H,N;\mathcal D)=
\sup\left\{R\geq0:\Pr_{Z\sim\mathcal D}
\left[I_\pi(R,H,N;Z)=1\right]\geq q\right\}.
\end{equation*}

The reliability target was 0.95. Each offer was conditional on the controller, workload distribution and call schedule. We selected the largest qualifying candidate on a finite development grid and tested it independently, using a pointwise probability lower bound.


This extension was specified after the earlier benchmark. One hundred development seeds (930000–930099) selected offers, and 300 previously unexamined confirmation seeds (960000–960299) tested them without reselection. The delivered-power recovery guard and feasible-side action rounding were corrected on development data before this controller version was confirmed. Protocol, source, controller and scenario hashes identify the evidence; original locked certificates remain separate historical results.


\subsection*{Community, jobs and event scenarios}

Community demand used the End-Use Load Profiles for the US Building Stock,\textsuperscript{32,33} combining 75\% detached-residence and 25\% small-office demand in the mixed 3A profile. These are building-stock simulations calibrated against measurements. Fifteen-minute power was averaged to hourly intervals and scaled to an 800-kW background peak. All five configurations prohibited exports and limited imports to 1,100 kW at the point of common coupling (PCC), the community's connection to the upstream grid. This rating accommodates the unchanged 576-GPU installation across the changing workload mixes; three initial development baselines exceeded the earlier 1,000-kW rating, so the common site rating was amended before formal PI optimisation or controller selection. Those baseline records were retained.


Job shapes came from the Alibaba Serverless Infrastructure execution summary,\textsuperscript{20} complementing an earlier public GPU trace.\textsuperscript{28} The source mix used all 40,522,321 released spans, weighted by requested GPU-equivalents times uncapped duration and retaining every workload category. A reproducible, class-balanced 100,000-record subset supplied low-priority training and offline-inference job shapes only. This summary lacked production deadlines and the original temporal correlations.


Each configuration contained 144 four-GPU nodes and offered 374.4 GPU-h per hour. Let s\_c denote a class share and a\_c its eligible fraction; total eligibility was \(f = \sum_c s_c a_c\). At f = 0.05, 0.10 and 0.20, a common permission multiplier was applied to the observed low-priority training and offline-inference resource-time shares, which together comprised 0.22028 of the source total. At f = 0.40, additional non-low-priority batch work was assumed to opt in. At f = 0.60, the source mix was explicitly changed by replacing 0.19058 of total work from online to offline inference. Online inference remained ineligible in every case. Flexible GPUs were round(576f), giving approximately 65\% mean utilisation in both pools; rigid utilisation used the actual integer remainder.


A common job template was generated once per seed. Class-specific GPU-hours were scaled across configurations while release times, deadlines, community values and the event anchor were checked for equality. Arrivals covered seven days, followed by a 48-h clearance tail. Training deadlines were sampled at 2–6 times runtime and clipped to 6–48 h; offline-inference deadlines were 1.5–4 times runtime and clipped to 2–24 h. The sampling specification and all five baseline-service checks are recorded in Supplementary Methods 1–3. These synthetic deadlines do not establish production service permissions.


\subsection*{Scheduling and power conversion}

Each job had a release time, class, GPU-hour requirement and deadline. Released work entered an earliest-deadline-first fluid queue. For class c in hour t,


\begin{equation*}
B_{c,t+1}=B_{c,t}+A_{c,t}-X_{c,t}-M_{c,t},
\end{equation*}

where B denotes backlog, A arrivals, X executed work and M unfinished work expiring at its deadline. Schedules respected hourly GPU capacity, work conservation, releases, deadlines, the allowed missed-work fraction and terminal backlog. Rigid and flexible work were represented separately. Work could be divided across hours; gang placement, non-pre-emptive execution and checkpoint latency were outside this abstraction. Each hourly step released arrivals before the control action, scheduled work, calculated power and advanced queue deadlines. Other time-step durations were rejected because the queue implementation advances one hourly bucket per step.


Data-centre power was calculated as


\begin{equation*}
\begin{aligned}
P^{\mathrm{DC}}_t=\frac{\mathrm{PUE}}{1000}\Big[&N_{\mathrm{node}}p_{\mathrm{node}}+N_{\mathrm{GPU}}p_{\mathrm{idle}}\\
&+N_{\mathrm{rigid}}u_{\mathrm{rigid}}(\bar p_{\mathrm{rigid}}-p_{\mathrm{idle}})
+\sum_c(p_c-p_{\mathrm{idle}})\frac{X_{c,t}}{\Delta t}\Big].
\end{aligned}
\end{equation*}

where power p is in watts, X is executed flexible GPU-hours and Δt = 1 h. Node overhead applies to all 144 nodes and idle power to all 576 GPUs. The rigid pool contains 518 GPUs at utilisation 0.65050193; its class-weighted active power is 291.42199 W/GPU. The reference flexible mix has 297.90848 W/GPU active power. At PUE 1.2 and idle power 13.935625 W/GPU, the operating peak is 51.840000 + 9.632304 + 112.202165 + 19.764511 = 193.438980 kW: node overhead, all-GPU idle draw, rigid increment and full flexible-pool reference increment, respectively. Thus the compact form \(P^{\mathrm{DC}}_t=P_{\mathrm{fixed}}+\sum_c e_cX_{c,t}\) uses \(P_{\mathrm{fixed}}=173.674469\) kW and \(e_c=\mathrm{PUE}(p_c-p_{\mathrm{idle}})/(1000\Delta t)\).


Calibration used four NVIDIA RTX PRO 6000 Blackwell Max-Q GPUs connected by PCIe without NVLink. Active training and offline-inference board power averaged 259.08 and 300.02 W/GPU, and idle power averaged 13.94 W/GPU. Two independent runs per active class supplied the fit and a third was held out. Simultaneous GPU observations were averaged within each run; independent runs were the statistical units.


The model used PUE 1.2 and 300 W/node fixed overhead. Online, development, other and unknown rigid work used a 300.02-W active-GPU engineering proxy, with separate 150- and 225-W controls. Rigid demand remained at its class-weighted mean. Board measurements calibrated the batch-work conversion; request latency and burst-level rigid power require application-level evidence.


Fixed-overhead sensitivity used 150, 300, 450 and 600 W/node (Supplementary Fig. 10; Supplementary Table 23). These change the operating peak to 167.52, 193.44, 219.36 and 245.28 kW. We added PUE × 144 × (p\_node − 300)/1000 to both power trajectories of all 600 independently confirmed reference schedules and rescored event and PCC constraints. This exact affine check tests feasibility of unchanged schedules, without selecting new offers or rerunning the controller. Baseline-relative reductions, backlog and waiting are invariant; absolute PCC headroom and the per-MW accounting conversion change.


\subsection*{Baselines, compute debt and successful delivery}

Every controlled episode had a no-demand-response baseline with the same community profile, jobs, deadlines and hardware. Compute debt measured the additional dynamic energy required to clear the controlled backlog relative to this baseline:


\begin{equation*}
D^{\mathrm{comp}}_t=\frac{\mathrm{PUE}}{1000}\sum_c\Delta B_{c,t}\left(p^{\mathrm{active}}_c-p^{\mathrm{idle}}\right).
\end{equation*}

where \(\Delta B_{c,t}\) is controlled minus baseline backlog in GPU-h, board powers are in W and compute debt is in kWh. Recovery and rebound were evaluated over the declared post-event window, with the clearance tail used to assess terminal backlog. For requested reduction R, delivered power was \(\Delta P_t=\max(0,P^{\mathrm{baseline}}_t-P^{\mathrm{control}}_t)\). Mean event delivery, \(\sum_{t\in E}\min(\Delta P_t,R)/(R|E|)\), had to reach 0.95, and every event hour also had to satisfy


\begin{equation*}
P^{\mathrm{control}}_t \leq P^{\mathrm{baseline}}_t-0.95R.
\end{equation*}

Success further required a deadline-miss fraction of eligible work no greater than 0.01, a rebound ratio no greater than 0.25, an event-plus-recovery peak-relief fraction of at least 0.50 and terminal backlog no greater than 0.02 of offered work. Rebound was the maximum positive controlled-minus-baseline PCC load in the 24-h recovery window, divided by peak event reduction. Mean and minimum interval delivery were reported separately. All criteria were fixed in the protocol; average delivered energy could not compensate for a failed hourly requirement.


The structural study additionally required zero total missed eligible work, allowing only 10\textsuperscript{−7} GPU-h for numerical round-off. Both response and no-response trajectories had to satisfy that zero criterion; all electrical and terminal-backlog thresholds were unchanged. The 0.1\% allowance was a secondary reported endpoint without offer selection. Baseline failures remained failures rather than exclusions. Earlier selected offers were also rescored at zero loss, explicitly as a retrospective check.


\subsection*{Relaxed planning and independent qualification}

The aggregate PI programme maximised R under cumulative class-release and class-deadline constraints. These are necessary aggregate conditions, but are a relaxation when job windows cross: work executed for an earlier, later-deadline job can satisfy a cumulative due-work inequality for a different job. Its retained optima therefore define relaxed planning envelopes, not generally job-feasible capacities. Their second order statistic among 100 scenarios is a 95\%/95\% tolerance statistic of that relaxed quantity, with achieved confidence 0.96292.\textsuperscript{34} It does not guarantee feasible work scheduling. Controller offers instead use completed queue simulations on independent seeds; their joint delivery-and-service probability is assessed by a one-sided 95\% Wilson lower bound, required to reach 0.95.\textsuperscript{35,36} The strict-window feasibility diagnostics use each work group’s own release and deadline.


The causal controller retained the six-hour robust MPC, historical-arrival uncertainty envelope and objective weights of the earlier implementation. It observed only released jobs, current queue state and the declared forecast. Its event cap required at least 95\% delivery. Because the rebound criterion divides recovery overshoot by actual peak delivered power, the corrected recovery cap used 0.25 × 0.95 × requested power rather than 0.25 × requested power. Continuous actions were rounded towards the feasible side when converted to float32. This electrical cap does not guarantee deadline feasibility. Candidates at 50\%, 75\%, 90\% and 100\% of the development Relaxed PI statistic were evaluated separately; no monotonicity between candidates was assumed.


\subsection*{Repeated commitments and structural controls}

All new structural tests used the corrected MPC wrapper, which retains a separate ceiling for each observed event-and-recovery window and clips the proposal to the strictest live ceiling. It uses elapsed baseline peaks and current observations, without future jobs or unannounced calls. Earlier comparisons with the original implementation and guarded greedy policy are retained in Supplementary Methods 6. New candidates were 25\%, 50\%, 75\% and 100\% of the fixed 5.898243186-kW reference request. Selection was separate for the 1\% and zero-miss endpoints and preceded independent confirmation; successful unselected comparators were not promoted.


At fixed 10\% work eligibility and 576 GPUs, offered utilisation was total work divided by installed GPU-hours, set to 50\%, 65\% or 80\%. Deadline slack was unchanged or halved with a one-hour floor; flexible GPU allocation was 10\%, with 20\% controls at 65\% and 80\% utilisation under tight deadlines. The final event started at hour 111 plus a seeded uniform integer from 0 to 23; preceding starts were spaced backwards by 16 or 24 h. All eight variants tested eight-hour calls, and the reference also tested four-hour calls at 16-h spacing. A fresh eight-hour final call at the same time removed the preceding calls. All recovery windows ended during 168 h of arrivals, followed by a 48-h clearance tail. Development seeds 984000–984099 and confirmation seeds 985000–985299 supplied 7,600 and 13,200 complete replays, respectively. Two pilot seeds were excluded.


For the external timing stress test, hourly submission counts from the Alibaba 2020 GPU trace\textsuperscript{28} supplied eight consecutive 168-h blocks. Counts set the hourly weights of synthetic class-specific work, preserving each class total and reference deadlines. The paired alternative permuted complete hourly blocks within the same week. Ten synthetic realisations per week crossed both orderings and 10\%/20\% GPU allocations with five frozen programme–request combinations, giving 1,600 replays without retuning. Results are reported by observed week and descriptively across 80 realisations per condition; they are not 80 independent production weeks.


\subsection*{Photovoltaic hosting and utilisation}

Community PCC power was


\begin{equation*}
P^{\mathrm{PCC}}_t=L_t+P^{\mathrm{DC}}_t+P^{\mathrm{ch}}_t-P^{\mathrm{dis}}_t-G^{\mathrm{PV}}_t,
\end{equation*}

with imports capped at 1,100 kW and exports prohibited. For each scenario we fixed the GPU installation and maximised PV nameplate subject to at most 5\% curtailment, zero missed GPU-hours and at most 2\% terminal backlog. BESS supplied 100-kW charge/discharge power and 200-kWh energy, with 0.95 charge and discharge efficiencies and initial and terminal state of charge of 50\%. A mixed-integer formulation prohibited simultaneous charging and discharging. All-scenario hosting was the minimum scenario capacity, reported only if every scenario was feasible.


The fixed-capacity programme used 500 kW of PV and lexicographically maximised local PV use, minimised grid import and then minimised battery throughput, retaining a 10\textsuperscript{−5}-kWh lock at each preceding optimum. Rigid and flexible schedules were paired on 100 inputs; mean utilisation gains used 10,000 paired bootstrap resamples (seed 20260911). No-BESS models are continuous linear programmes, so a relative MIP gap does not set their precision. We repeated all 800 primary 10\% PV optimisations and 200 no-BESS GPU20 hosting optimisations with HiGHS\textsuperscript{37} relative and absolute MIP gaps 10\textsuperscript{−7} and primal, dual and MIP feasibility tolerances 10\textsuperscript{−8}. Each solve used one thread and a 120-s limit. All fixed-PV programmes resolved; 34 flexible BESS hosting problems reached the time limit. Their feasible and dual bounds, retained in full, did not affect the minimum hosting boundary over 100 scenarios. Primary-objective bounds separate numerical uncertainty from paired sampling intervals (Supplementary Table 24).


\subsection*{Economic participation}

Economic screening used each configuration's development-selected, independently tested four- or eight-hour zero-notice offer. Controlled and no-response trajectories supplied paired energy, delay, deferred work and service outcomes from the first call through the 48-h clearance tail. Performance revenue used non-negative PCC reduction capped at the offer in each event hour. Annual net value retained the following accounting structure; inactive or unmeasured terms were stated explicitly rather than inferred from the hardware measurements:


\begin{equation*}
V_{\mathrm{annual}}(K)=R_{\mathrm{capacity}}+R_{\mathrm{performance}}+V_{\mathrm{energy}}+V_{\mathrm{connection}}-C_{\mathrm{enable}}-C_{\mathrm{reserve}}-C_{\mathrm{opportunity}}-C_{\mathrm{delay}}-C_{\mathrm{SLA}}-C_{\mathrm{checkpoint}}-C_{\mathrm{penalty}}.
\end{equation*}

Reference prices were US\$50 per MWh of capped delivered energy, US\$0.10 per kWh of incremental electricity and US\$0.005 per GPU-h per hour of additional waiting. We solved for the capacity payment making the fifth percentile of 2,000 bootstrap annual net values non-negative. This was a decision statistic, not a confidence interval. Natural-slack participation incurred no reserve or displacement charge. The reserve case annualised 0.15 abstract reserved PCC-side kW per offered kW at US\$2,500 per reserved kW, four-year life, 8\% discount, 10\% salvage and 3\% annual operation and maintenance. The displacement case added an assumed US\$0.25 per deferred GPU-h; this exposure was distinct from delay cost and was not a measured revenue loss. Connection benefits and checkpoint costs were unpriced; the reference missed-work price was zero. The model rescheduled work on the same powered GPUs and assumed no DR-induced change in useful life. Accordingly, no incremental hardware-degradation charge was assigned; the power calibration did not test this lifetime assumption, and existing-fleet depreciation was not charged again as an incremental response expense.


Enablement costs combined a site-wide annual charge with a variable per-kilowatt charge:


\begin{equation*}
C_{\mathrm{enable}}(K)=\mathbf 1[K>0]C_{\mathrm{fixed,site}}+c_{\mathrm{variable}}K.
\end{equation*}

The reference fixed site charge was US\$25,000 per year and the variable enablement charge was zero. Accounting multiplied offered power and physical ledger quantities by 1000 divided by each configuration’s operating peak, then applied the site charge once. Annual draws sampled 50 independent single events or 12 independent complete four-call series, retaining failures. The external-source register\textsuperscript{38–43} contextualises the financial sensitivity ranges; Scope and interpretation defines the limits of this proportional accounting.


Repeated-programme costs were accumulated once per hour over the complete paired series. A programme that did not qualify was identified as such; its conditional cost did not create an eligible offer. Prices of US\$0, 0.25 and 1 per missed GPU-h were separately evaluated, and contract-specific non-performance penalties were not priced. Independent annual draws do not reproduce a continuously operating year.


The fee decomposition used common annual draws and the same linear-interpolation weights at the total net-cost 95th percentile for every component; it did not add marginal component quantiles. Additional screens crossed site fees US\$0, 1,000, 2,500, 5,000, 10,000 and 25,000, waiting prices US\$0, 0.005 and 0.01 per GPU-h/h, and missed-work prices US\$0 and 1 per GPU-h. Sharing a US\$25,000 fee among ten or five resources gave US\$2,500 or 5,000 without a diversification benefit. For qualified product j with accounting capacity K\_j and annual net operating cost O\_j at its 95th percentile, the fifth-percentile net value was P\_j K\_j − O\_j − F. Duration-price boundaries compared these values and the zero value of not participating. They are algebraic consequences of the stated ledgers and prices, not estimated market tariffs.


\subsection*{Sensitivity analysis and statistical reporting}

The five eligibility scenarios used a declared matching rule for GPU allocation and therefore are not a one-factor experiment across the whole range. Two additional controls kept 10\% eligibility and the same jobs while assigning 20\% or 30\% of GPUs to the flexible pool. Two further controls changed only the unmeasured rigid-class active-power proxy to 150 or 225 W. Each control repeated PI planning, the same development-selected fixed requests, PV analysis and economic accounting. Confirmation never retuned the controller or selected new offers for these controls. Earlier broader parameter, regional and joint-share studies are retained with their original workload scope in the archived v0.19 Supplementary Information.


Economic sensitivities recomputed payments from the fixed paired ledgers with common annual draws, without changing physical trajectories or offer selection. They crossed delay prices, site costs and assumed displaced-compute values, and separately varied missed-work prices. Hardware procurement, actual online-service costs and a continuous-year market simulation were not inferred from these screens.


Independent simulated scenarios were the statistical units for model qualification; calls within a scenario were scored jointly. Relaxed PI tolerance statistics, Wilson probability lower bounds and paired bootstrap intervals were pointwise. All 20 notice comparisons had zero discordant success pairs, so we report their outcome identity directly. Hardware intervals retained the independent-run analysis. Source Data contain individual outcomes, frozen selections, complete paired ledgers and plotting summaries, with hashes linking source, protocol and controller versions.


Structural contrasts resampled 300 complete paired scenario seeds 5,000 times. Exact McNemar tests were Holm-adjusted over 36 binary structural contrasts; the 16 additional matched-target electrical contrasts formed a separate family. Both service endpoints used the pointwise one-sided 95\% Wilson lower-bound criterion of 0.95. External workload summaries retained the observed week as their sampling unit; conditional empirical-mixture qualification is defined separately below.


Local refinement fixed the reference workload and controller and tested fractions 0.50–0.75 for H8P16 and 0.75–1.00 for H4P16, in steps of 0.05 of the unchanged 5.898243186-kW request. H denotes duration and P start spacing. Seeds 988000–988099 supplied development; endpoint-specific selections were frozen before seeds 989000–989299 supplied confirmation. Only selected candidates and prespecified coarse comparators entered confirmation, with no reselection. The grid ceiling for eight hours was selected under both standards, so the test did not bracket a continuous maximum. Earlier seeds were not pooled with these data.


Perfect-information diagnostics fixed each request and its four event clocks, baseline and 24-h recovery windows. Non-negative execution variables existed only between each class/release/deadline work group’s own release and deadline; execution, misses and terminal work conserved every group. Constraints retained GPU and PCC limits, 95\% hourly and capped-mean delivery, 25\% rebound relative to actual peak event reduction, 50\% full-window peak relief and the original terminal allowance. Binary peak selectors represented the rebound denominator exactly. HiGHS solved the resulting mixed-integer feasibility problem. Thirty-three prespecified structural cases comprised five configurations, three previously used seeds and two service standards, plus three reference 4.42-kW checks; five affected external weeks supplied an additional first-shortage diagnostic each. These mechanism-selected cases were not prevalence samples. All 11 feasible witnesses passed independent work-group and electrical checks. A 90-s unresolved solve was retained and resolved as infeasible after extending only its time limit to 600 s. Full-information feasibility was not treated as causal realizability.


The resource-time transfer test reconstructed each processed job’s GPU-seconds as the sum of requested GPU-equivalents × task launch-to-completion duration across its terminated positive-resource tasks.\textsuperscript{28} All 714,903 processed jobs matched the raw task reconstruction; multiplying total requested GPUs by job submission-to-completion time instead differed for 131,379 jobs. The test retained trace hours 168–1511, whose processed time origin is 542,323 s after the raw origin. Within each week, either job counts or submitted task resource-time set hourly weights; each class’s weekly work total was normalised to the reference total. Identical synthetic job classes, deadlines, community traces, event clocks and paired whole-hour permutations were used at each seed (990000 + 100w + r; eight weeks × ten realisations). Fixed 2.95- and 4.42-kW requests gave 640 complete replays. All 320 no-response input variants passed zero-loss service; no failed variant was excluded. Aggregate counts are descriptive. Allocated resource-time is not a sensor measurement of busy GPU time, and the test retains divisible work and synthetic service rules.


The follow-up capacity experiment conditioned on those eight fixed weekly profiles. For each new seed, an independent uniform draw selected one week; fresh job templates, community conditions and a random final-call phase were paired between count and resource-time weights. Four eight-hour calls started 16 h apart with zero notice. Development seeds 1010000–1010099 tested fractions 1/64, 1/32, 1/16, 1/8, 1/4, 3/8, 1/2, 3/4 and 1 of 5.898243186 kW. The largest candidate with a one-sided 95\% Wilson lower bound ≥0.95 was frozen separately for each weight and service endpoint. Confirmation seeds 1011000–1011299 then tested the selected offers and a prespecified 1/2-request comparator, retaining baseline failures. The 1,800 development and 900 confirmation replays did not reuse the earlier 80 realisations for selection. Inference concerns simulated draws from the fixed eight-week empirical mixture; it does not treat them as newly observed production weeks.


\subsection*{Scope and interpretation}

Work eligibility is permission over offered GPU-hours; operating peak is the configured reference-mix facility denominator. A selected, confirmed offer is a finite-grid policy result for its stated distribution. Exact-window witnesses diagnose individual cases, whereas probabilistic qualification retains all sampled cases. Resource-time weights describe requested task allocation and duration; the count-only control deliberately removes that information and is not attributed to Caprara et al. or another published scheduler. Proportional 1-MW accounting rescales the model's ledgers and applies one site fee; no additional facility, geographic aggregation or diversification is simulated. Economic comparisons retain the workload distribution and products on which their offers were confirmed.


\section*{Data Availability}

Community demand was obtained from the NREL End-Use Load Profiles through the OEDI building-stock data lake.\textsuperscript{32,33} Workload records came from the official Alibaba cluster-trace-gpu-v2026 release,\textsuperscript{20} with the 2020 GPU trace supplying the separate chronological submission test.\textsuperscript{28} Retrieval locations, preprocessing records and original input hashes are listed in data/manifests/sources.yaml. Raw third-party releases are accessed from their original providers. The accompanying Source Data bundle contains the full class-composition audit, scenario and controller identities, development selection, all trial-level outcomes, complete development and confirmation trajectories and paired economic ledgers, renewable optimisation results, calibration inputs and ready-to-plot tables. CSV and Parquet files have data dictionaries and SHA-256 manifests. Earlier benchmark outputs are preserved in a separately labelled archive. Public archival deposition of this revised bundle remains pending at [ZENODO DOI TO BE ADDED]. Submission counts, processed submission-time records, frozen scenarios and all 22,400 new replays accompany the structural Source Data. The focused extension adds 3,040 full replays, 656,640 hourly rows, 38 resolved exact-window feasibility diagnoses and task-level resource-time reconstruction records.


\section*{Code Availability}

The AIDRBench project is hosted at \url{https://github.com/daxiang415/AIDRBench}. The accompanying revision bundle includes the exact study drivers, corrected-controller implementation, frozen protocols, analysis scripts and standalone Python redraw commands used here. Public archival release of the revised submission code and its software licence remain pending author confirmation; no new release DOI is claimed.


\section*{References}
{\small
\begin{enumerate}[label=\arabic*.,leftmargin=*,itemsep=2.5pt]
\item International Energy Agency. \emph{Energy and AI} (IEA, 2025). \url{https://www.iea.org/reports/energy-and-ai}
\item Lin, L. et al. Exploding AI power use: an opportunity to rethink grid planning and management. In \emph{Proceedings of the 15th ACM International Conference on Future and Sustainable Energy Systems} 434–441 (ACM, 2024). \url{https://doi.org/10.1145/3632775.3661959}
\item Colangelo, P. et al. AI data centres as grid-interactive assets. \emph{Nat. Energy} \textbf{11}, 254–261 (2026). \url{https://doi.org/10.1038/s41560-025-01927-1}
\item Wierman, A., Liu, Z., Liu, I. \& Mohsenian-Rad, H. Opportunities and challenges for data center demand response. In \emph{2014 International Green Computing Conference} 1–10 (IEEE, 2014). \url{https://doi.org/10.1109/IGCC.2014.7039172}
\item Li, J., Bao, Z. \& Li, Z. Modeling demand response capability by Internet data centers processing batch computing jobs. \emph{IEEE Trans. Smart Grid} \textbf{6}, 737–747 (2015). \url{https://doi.org/10.1109/TSG.2014.2363583}
\item Zhang, Y., Wilson, D. C., Paschalidis, I. Ch. \& Coskun, A. K. HPC data center participation in demand response: an adaptive policy with QoS assurance. \emph{IEEE Trans. Sustain. Comput.} \textbf{7}, 157–171 (2022). \url{https://doi.org/10.1109/TSUSC.2021.3077254}
\item Zhang, Y., Wilson, D. C., Paschalidis, I. Ch. \& Coskun, A. K. A data center demand response policy for real-world workload scenarios in HPC. In \emph{2021 Design, Automation \& Test in Europe Conference \& Exhibition} 282–287 (IEEE, 2021). \url{https://doi.org/10.23919/DATE51398.2021.9474075}
\item Liu, Z., Wierman, A., Chen, Y., Razon, B. \& Chen, N. Data center demand response: avoiding the coincident peak via workload shifting and local generation. \emph{Perform. Evaluation} \textbf{70}, 770–791 (2013). \url{https://doi.org/10.1016/j.peva.2013.08.014}
\item Radovanović, A. et al. Carbon-aware computing for datacenters. \emph{IEEE Trans. Power Syst.} \textbf{38}, 1270–1280 (2023). \url{https://doi.org/10.1109/TPWRS.2022.3173250}
\item Riepin, I., Brown, T. \& Zavala, V. M. Spatio-temporal load shifting for truly clean computing. \emph{Adv. Appl. Energy} \textbf{17}, 100202 (2025). \url{https://doi.org/10.1016/j.adapen.2024.100202}
\item Zheng, J., Chien, A. A. \& Suh, S. Mitigating curtailment and carbon emissions through load migration between data centers. \emph{Joule} \textbf{4}, 2208–2222 (2020). \url{https://doi.org/10.1016/j.joule.2020.08.001}
\item Goiri, Í., Katsak, W., Le, K., Nguyen, T. D. \& Bianchini, R. Parasol and GreenSwitch: managing datacenters powered by renewable energy. In \emph{Proceedings of the Eighteenth International Conference on Architectural Support for Programming Languages and Operating Systems} 51–64 (ACM, 2013). \url{https://doi.org/10.1145/2451116.2451123}
\item Lechowicz, A. et al. The online pause and resume problem: optimal algorithms and an application to carbon-aware load shifting. \emph{Proc. ACM Meas. Anal. Comput. Syst.} \textbf{7}, 1–32 (2023). \url{https://doi.org/10.1145/3626776}
\item Zhao, D. et al. Sustainable supercomputing for AI: GPU power capping at HPC scale. In \emph{Proceedings of the 2023 ACM Symposium on Cloud Computing} 588–596 (ACM, 2023). \url{https://doi.org/10.1145/3620678.3624793}
\item Caprara, A., Yu, Y., Teng, F., Junyent-Ferré, A., Bullich-Massagué, E. \& Aragüés-Peñalba, M. Data center workload flexibility for power system demand response: evidence from Alibaba traces. \emph{Int. J. Electr. Power Energy Syst.} \textbf{178}, 111940 (2026). \url{https://doi.org/10.1016/j.ijepes.2026.111940}
\item Chen, Y. \& Zheng, X. To defer or to shift? The role of AI data center flexibility on grid interconnection. In \emph{Proceedings of the 2026 ACM Sustainability Week} 322–327 (ACM, 2026). \url{https://doi.org/10.1145/3765611.3815593}
\item Zhao, M., Wang, X. \& Mo, J. Workload and energy management of geo-distributed datacenters considering demand response programs. \emph{Sustain. Energy Technol. Assess.} \textbf{55}, 102851 (2023). \url{https://doi.org/10.1016/j.seta.2022.102851}
\item Wiesner, P., Behnke, I., Scheinert, D., Gontarska, K. \& Thamsen, L. Let's wait awhile: how temporal workload shifting can reduce carbon emissions in the cloud. In \emph{Proceedings of the 22nd International Middleware Conference} 260–272 (ACM, 2021). \url{https://doi.org/10.1145/3464298.3493399}
\item Li, M. Beyond Scalar Flexibility: From Eligible AI Workloads to Dependable Load Relief. Preprint at \url{https://arxiv.org/abs/2609.05406} (2026).
\item Li, S. et al. Heterogeneity at hyperscale: characterization and scheduling of large production AI clusters at Alibaba. In \emph{20th USENIX Symposium on Operating Systems Design and Implementation} 2187–2203 (USENIX Association, 2026).
\item Cioara, T. et al. Optimized flexibility management enacting data centres participation in smart demand response programs. \emph{Future Gener. Comput. Syst.} \textbf{78}, 330–342 (2018). \url{https://doi.org/10.1016/j.future.2016.05.010}
\item Wang, W., Abdolrashidi, A., Yu, N. \& Wong, D. Frequency regulation service provision in data center with computational flexibility. \emph{Appl. Energy} \textbf{251}, 113304 (2019). \url{https://doi.org/10.1016/j.apenergy.2019.05.107}
\item Su, T. et al. Grid-enhancing technologies for clean energy systems. \emph{Nat. Rev. Clean Technol.} \textbf{1}, 16–31 (2025). \url{https://doi.org/10.1038/s44359-024-00001-5}
\item Loji, K., Sharma, S., Sharma, G. \& Rawat, T. Multiobjective distribution system operation with demand response to optimize solar hosting capacity, voltage deviation index and network loss. \emph{Sci. Rep.} \textbf{15}, 300 (2025). \url{https://doi.org/10.1038/s41598-024-82379-7}
\item Fu, Y., Bai, H., Cai, Y., Yang, W. \& Li, Y. Optimal configuration method of demand-side flexible resources for enhancing renewable energy integration. \emph{Sci. Rep.} \textbf{14}, 7658 (2024). \url{https://doi.org/10.1038/s41598-024-58266-6}
\item Davies, D. M. et al. Combined economic and technological evaluation of battery energy storage for grid applications. \emph{Nat. Energy} \textbf{4}, 42–50 (2019). \url{https://doi.org/10.1038/s41560-018-0290-1}
\item Li, K. et al. Facilitating megacity electricity decarbonization via grid-interactive demand-side resource management. \emph{Nat. Commun.} (2026). \url{https://doi.org/10.1038/s41467-026-76799-4}
\item Weng, Q. et al. MLaaS in the wild: workload analysis and scheduling in large-scale heterogeneous GPU clusters. In \emph{19th USENIX Symposium on Networked Systems Design and Implementation} 945–960 (USENIX Association, 2022).
\item PJM Interconnection. \emph{Manual 18: PJM Capacity Market}, sections 4.3.2 and 8 (accessed 11 September 2026). \url{https://www.pjm.com/-/media/DotCom/documents/manuals/m18}
\item Zhou, Y., Mancarella, P. \& Mutale, J. Framework for capacity credit assessment of electrical energy storage and demand response. \emph{IET Gener. Transm. Distrib.} \textbf{10}, 2267–2276 (2016). \url{https://doi.org/10.1049/iet-gtd.2015.0458}
\item Feng, J. et al. Evaluating demand response impacts on capacity credit of renewable distributed generation in smart distribution systems. \emph{IEEE Access} \textbf{6}, 14307–14317 (2018). \url{https://doi.org/10.1109/ACCESS.2017.2745198}
\item Wilson, E. et al. \emph{End-Use Load Profiles for the U.S. Building Stock: Methodology and Results of Model Calibration, Validation, and Uncertainty Quantification}. NREL/TP-5500-80889 (National Renewable Energy Laboratory, 2022). \url{https://doi.org/10.2172/1854582}
\item Parker, A. et al. \emph{ComStock Reference Documentation: Version 1}. NREL/TP-5500-83819 (National Renewable Energy Laboratory, 2023). \url{https://doi.org/10.2172/1967948}
\item Vangel, M. G. One-sided nonparametric tolerance limits. \emph{Commun. Stat. Simul. Comput.} \textbf{23}, 1137–1154 (1994). \url{https://doi.org/10.1080/03610919408813222}
\item Wilson, E. B. Probable inference, the law of succession, and statistical inference. \emph{J. Am. Stat. Assoc.} \textbf{22}, 209–212 (1927). \url{https://doi.org/10.1080/01621459.1927.10502953}
\item Brown, L. D., Cai, T. T. \& DasGupta, A. Interval estimation for a binomial proportion. \emph{Stat. Sci.} \textbf{16}, 101–133 (2001). \url{https://doi.org/10.1214/ss/1009213286}
\item Huangfu, Q. \& Hall, J. A. J. Parallelizing the dual revised simplex method. \emph{Math. Program. Comput.} \textbf{10}, 119–142 (2018). \url{https://doi.org/10.1007/s12532-017-0130-5}
\item CoreWeave, Inc. \emph{Annual Report (Form 10-K) for the Year Ended 31 December 2025}. US Securities and Exchange Commission (2026). \url{https://www.sec.gov/Archives/edgar/data/1769628/000176962826000104/crwv-20251231.htm}
\item Amazon.com, Inc. \emph{Annual Report (Form 10-K) for the Year Ended 31 December 2025}. US Securities and Exchange Commission (2026). \url{https://www.sec.gov/Archives/edgar/data/1018724/000101872426000004/amzn-20251231.htm}
\item Microsoft Corporation. \emph{Annual Report (Form 10-K) for the Fiscal Year Ended 30 June 2025}. US Securities and Exchange Commission (2025). \url{https://www.sec.gov/Archives/edgar/data/789019/000095017025100235/msft-20250630.htm}
\item National Renewable Energy Laboratory. \emph{2024 Annual Technology Baseline: Financial Cases and Methods} (NREL, 2024). \url{https://atb.nrel.gov/electricity/2024/financial_cases_\%26_methods}
\item Lawrence Berkeley National Laboratory. \emph{Demand Response Advanced Controls Framework and Cost Assessment} (LBNL, 2017). \url{https://eta-publications.lbl.gov/sites/default/files/demand_response_advanced_controls_framework_and_cost_assessment_final_published.pdf}
\item PJM Interconnection. \emph{2027/2028 Base Residual Auction Report} (PJM, 2025). \url{https://www.pjm.com/-/media/DotCom/markets-ops/rpm/rpm-auction-info/2027-2028/2027-2028-bra-report.pdf}
\end{enumerate}}
\section*{Acknowledgements}

[AUTHOR INPUT NEEDED: funding, facilities and non-author contributions.]


\section*{Author Contributions}

[AUTHOR INPUT NEEDED: CRediT-aligned author contributions.]


\section*{Competing Interests}

[AUTHOR INPUT NEEDED: competing-interests declaration.]


\clearpage
\section*{Figure Legends}
\small\setlength{\parskip}{3pt}

\begin{minipage}{\textwidth}
\subsection*{Figure 1 | AI workload composition, deferred scheduling and assessment of demand-response commitments}

\textbf{a,} Conceptual overview linking workload representation, job-constrained scheduling and grid commitments. The task bars illustrate postponement and recovery without a numerical time scale. Online service remains immediate; operators decide which training and offline-inference work may wait. \textbf{b,} Requested GPU-hour shares from all 40,522,321 released execution spans, using uncapped duration and including all six workload classes. These descriptive source totals characterise computing work rather than electricity consumption. \textbf{c,} Evidence-to-decision sequence: reject requests exceeding supply or strict task-window feasibility, investigate control where a feasible full-information schedule exists, qualify the causal policy independently and then compare participation costs. Qualification retains successful and failed trials. An unresolved feasibility solve supplies no definitive decision. Source Data retain full-precision class and scenario definitions.
\end{minipage}\par\vspace{7pt}
\begin{minipage}{\textwidth}
\subsection*{Figure 2 | Single-event response capacity and reliability across workload eligibility, event duration and advance notice}

\textbf{a,} Relaxed perfect-information (PI) tolerance statistics from 100 confirmation scenarios (95\% reliability and at least 95\% confidence; horizontal marks) and development-selected, independently confirmed single-event offers (circles, four hours; squares, eight hours). Their different constructions mean that their ordering does not imply an information advantage. \textbf{b,} The same offers divided by each configuration's operating peak. Four- and eight-hour markers coincide because their selected capacities are identical. \textbf{c,} A compact summary of all five eligibility settings and three notice periods: paired success flags are identical at 0-, 2- and 6-h notice, with 295/300 four-hour and 292/300 eight-hour successes in each setting. Values beside the counts are one-sided 95\% Wilson probability lower bounds, evaluated separately for each condition. The target is 95\% success. Orange and teal identify four- and eight-hour events in this figure. The PI comparison uses aggregate release/deadline relaxations and does not certify individual-task feasibility. Additional business assumptions at 40\% (*) and 60\% (†) are detailed in Supplementary Table 2.
\end{minipage}\par\vspace{7pt}
\begin{minipage}{\textwidth}
\subsection*{Figure 3 | Repeated-response offers and delivery outcomes under submission-count weighting, task-resource-time weighting and hourly reordering}

\textbf{a,} Selected eight-hour repeated offers confirmed on 300 independent draws from an equal mixture of eight fixed observed weeks. Both service standards pass in 300/300 draws for each selected offer (one-sided 95\% Wilson lower bound 0.9911). Bar lengths encode capacity; the twofold ratio compares finite-grid selections. \textbf{b,} All nine development requests per workload representation, each evaluated on 100 model draws before confirmation. Lines join tested points; circles and solid lines allow at most 1\% missed work, whereas squares and dashed lines require zero missed work. The horizontal line marks the 95\% qualification target. \textbf{c,} Zero-miss diagnoses at 5.90 kW for three prespecified seeds per configuration, distinguishing immediate shortage, other strict-window infeasibility, PI feasibility with causal failure and feasibility under both tests. \textbf{d,} Paired weight/order comparisons at 2.95 kW under both service standards. \textbf{e,} Original-order zero-miss successes by observed week; each cell retains all ten synthetic realisations. Panels d,e report descriptive counts. Weekly class totals and service rules are paired, all baselines meet zero-miss service, and all 49 original-order resource-time failures encounter shortage. Permutation changes hourly order and alignment with calls and community demand. Reference-distribution duration products are assessed separately in Supplementary Fig. 9a.
\end{minipage}\par\vspace{7pt}
\begin{minipage}{\textwidth}
\subsection*{Figure 4 | Available-load shortages during repeated calls and power and backlog trajectories through recovery}

\textbf{a,} Paired minimum event-hour supply margins for all 80 task-resource-time realisations at 2.95 kW, comparing original and permuted orders. Margin is baseline PCC power minus community and idle-plus-rigid data-centre demand minus 95\% of the request, minimised across all calls. Zero lines separate shortage and no-shortage cases; quadrant counts are descriptive. \textbf{b,} Available baseline reduction under chronological count and resource-time weighting for protocol seed 990000. The horizontal line is the 95\% delivery requirement and shading marks four calls. All 216 h are retained in Source Data; the panel displays the 168-h arrival horizon. \textbf{c,d,} Aligned power difference from baseline and additional backlog over all 216 h for reference seed 989000, at qualified four- and eight-hour offers of 5.60 and 4.42 kW. Both products make four calls starting 16 h apart. Numbered shading shows the eight-hour windows; four-hour calls share their starts. Solid lines show operation and dotted lines show requests. Negative power reduction indicates recovery above baseline. The grey 168–216-h tail permits clearance; both illustrated series finish with zero missed work. Seeds were selected by identifier, without inspecting outcomes. Panels a,b and c,d use different workload distributions. These examples explain operation; independent reliability evidence is in Supplementary Fig. 9a.
\end{minipage}\par\vspace{7pt}
\begin{minipage}{\textwidth}
\subsection*{Figure 5 | Response planning, participation costs and solar hosting across GPU allocations and rigid-load power assumptions}

\textbf{a,b,} Four- and eight-hour relaxed-PI tolerance statistics across flexible GPU allocations or rigid-class active-power proxies, retaining 10\% eligibility, 576 GPUs and paired jobs. Each statistic uses 100 confirmation scenarios and aggregate release/deadline constraints; it does not guarantee individual-task feasibility. The fixed 5.90-kW causal request has 295/300 four-hour and 292/300 eight-hour successes across these controls (Supplementary Fig. 4c). Labels in b report operating peaks. \textbf{c,d,} Conditional four-hour single-event participation thresholds at the unchanged offer, from all 300 confirmation ledgers, 50 independent calls per annual draw, a US\$25,000 annual site fee and proportional 1-MW accounting. Points are 95th-percentile annual-net-cost thresholds. \textbf{e,} Flexible minus rigid PV hosting capacity as GPU allocation changes; each hosting capacity is the minimum across 100 paired full-information, zero-miss scenarios. These are differences of minima, rather than uncertainty intervals. Lines join evaluated settings. The 10\% primary points and no-storage 20\% allocation point use the strict repeat audit; remaining PV points retain their original settings (Supplementary Table 24). Rigid power proxies are engineering assumptions, not measurements of online inference.
\end{minipage}\par\vspace{7pt}
\begin{minipage}{\textwidth}
\subsection*{Figure 6 | Additional waiting, participation costs and service choice for qualified four- and eight-hour repeated responses}

\textbf{a,} Empirical cumulative distributions of additional waiting exposure per four-call series for the two qualified reference products, using all 300 confirmation sequences per product, including failures. Dotted vertical lines mark means. No smoothing or density fitting is applied. \textbf{b,} Operating cost components at the total annual 95th-percentile rank, with zero site fee and US\$0.005/GPU-h/h waiting; diamonds mark the net threshold. \textbf{c,} Participation thresholds at the three tested waiting prices. \textbf{d,} The product with the largest fifth-percentile annual net value, including non-participation, as the two capacity prices vary at US\$0.005/GPU-h/h waiting and zero site fee. The dotted diagonal marks equal capacity prices. Region boundaries follow the frozen cost and capacity accounting, with ties along their shared edges. All panels concern reference products qualified at a 95\% success-probability target and a zero-miss success criterion: 5.60 kW for four hours and 4.42 kW for eight hours, with 16-h start spacing. Monetary panels use proportional 1-MW accounting, twelve independent four-call series per year, US\$0.10/kWh electricity, US\$50/MWh capped delivered-energy revenue, zero missed-work price and 2,000 annual draws retaining failures. Prices remain conditional on this reference workload; they do not price the separate resource-time offers.
\end{minipage}\par\vspace{7pt}
\end{document}
